% ID: SEL-D0030-OA-E01
% ESTADO: validada. S031, cotejada con examen y criterios.
\textbf{a)} Sean $x,y$ los kilogramos diarios de maíz y pienso
por oveja. Los mínimos nutricionales y la no negatividad exigen
\[
600x+300y\ge1800,\qquad 200x+600y\ge2400,\qquad x,y\ge0.
\]
La función que se debe minimizar es $C(x,y)=0.50x+0.25y$,
en euros diarios por oveja.

\textbf{b)} Las cotas verticales son $x/2-1\le y\le5-x$.
En $x=0$ van de $-1$ a $5$. Al igualarlas resulta $x=4$ e $y=1$.
Por tanto los vértices son $(0,-1)$, $(0,5)$ y $(4,1)$.
\begin{center}
\setlength{\unitlength}{7mm}
\begin{picture}(9,8)(-2,-2)
\qbezier(0,-1.6)(0,2)(0,5.6)\put(0,5.7){$y$}
\qbezier(0,0)(2.5,0)(5,0)\put(5.2,0){$x$}
\qbezier(0,-1)(2,0)(4,1)
\qbezier(4,1)(2,3)(0,5)
\put(-0.2,5){\makebox(0,0)[r]{$(0,5)$}}
\put(-0.2,-1){\makebox(0,0)[r]{$(0,-1)$}}
\put(4.2,1){$(4,1)$}
\end{picture}
\end{center}
La función lineal vale $-3$, $15$ y $19$ en esos vértices.
El máximo es $\boxed{19}$ y se alcanza sólo en $(4,1)$.
